\documentclass[journal]{IEEEtran}

\usepackage{amsmath}
\usepackage{booktabs}
\usepackage{cite}

\title{LLM-Enhanced Explainable Fault Diagnosis Toolkit: An Evidence-Grounded Draft for IEEE Transactions Review}

\author{Li~Qi%
\thanks{This draft is a submission-readiness checkpoint for the UXFD paper family. Quantitative performance, human-study, TOP-baseline, and SOTA claims are intentionally withheld until accepted run artifacts are available.}}

\begin{document}

\maketitle

\begin{abstract}
Industrial fault diagnosis systems increasingly combine signal-processing models
with natural-language interfaces. This paper draft describes an
LLM-enhanced explainable fault diagnosis toolkit that converts structured
diagnostic evidence into auditable text and dialogue responses. The current
IEEE package defines the intended architecture, evidence contract, baseline
surface, and reviewer-facing limitations. It does not claim submission-ready
performance: accepted multi-seed experiments, run metadata, latency metrics,
unsupported-claim rates, TOP representative artifacts, and GPU provenance are
still required before any SOTA or human-centered decision-support claim can be
made.
\end{abstract}

\begin{IEEEkeywords}
Explainable fault diagnosis, large language models, evidence grounding,
industrial AI, human-centered decision support.
\end{IEEEkeywords}

\section{Introduction}
\IEEEPARstart{I}{ndustrial} diagnostic models can classify fault states from
vibration and other condition-monitoring signals, but their outputs often remain
difficult for engineers to inspect, question, and reuse in maintenance
decisions. A language interface can reduce this gap only if generated text is
grounded in structured diagnostic evidence. Ungrounded prose is unsafe for
industrial diagnosis because it may overstate model confidence, invent causal
explanations, or hide uncertainty.
This evidence-first position is consistent with established local explanation
work such as LIME and SHAP~\cite{ribeiro2016why,lundberg2017unified}, and with
the need to distinguish attention-style representations from verified
diagnostic evidence~\cite{vaswani2017attention}.

This draft frames the LLM-enhanced explainable fault diagnosis toolkit as an
evidence-grounded reporting layer. The toolkit consumes structured model
outputs, maps them into an intermediate representation, and generates natural
language responses whose claims must be traceable to signal statistics,
diagnostic predictions, model explanations, or a controlled knowledge base. The
draft is deliberately conservative: all numerical performance and user-study
claims from earlier Markdown drafts are excluded until accepted artifacts exist.

\section{Method Overview}
\subsection{Structured Evidence Interface}
The toolkit is designed around an intermediate representation containing fault
type, confidence, signal statistics, frequency features, processing-path
information, and optional operator or attribution evidence. This representation
is the only allowed source for generated diagnostic statements. The evidence
interface supports transparent signal-processing models, post-hoc feature
summaries, and rule-based explanation modules.

\subsection{Local Template and LLM Backends}
The current package includes a deterministic local template backend for
offline smoke testing. External LLM providers can be integrated through the
provider manager, but accepted paper evidence must record provider settings,
prompt versions, failure modes, and whether the generated text is exact LLM
output or a local proxy. This separation is necessary for fair baseline
comparison and reproducibility, especially when comparing against recent
time-series foundation or LLM adaptation methods such as Time-LLM and
MOMENT~\cite{jin2024timellm,goswami2024moment}.

\subsection{Dialogue and Safety Constraints}
The interactive interface supports common diagnostic intents such as cause
analysis, severity assessment, maintenance guidance, monitoring advice, and
technical-detail queries. Each response must preserve the fault label,
confidence, and evidence scope from the current diagnostic context. Unsupported
claims, missing evidence fields, and unavailable retrieval context must be
reported as limitations rather than silently filled by generated text.

\section{Submission Evidence Protocol}
Accepted evidence must live under a paper-local
\texttt{results/llm\_evidence/} directory. Each run directory must include a
\texttt{run\_meta.yaml} file and a \texttt{metrics.json} file. The metadata
must record the command, working directory, seed, prompt set, model or template
backend, precision or quantization mode, runtime, artifact paths, and local
device provenance. If a run is blocked by hardware, data, dependency, or policy
constraints, the failure reason must be recorded explicitly.

The required metrics include diagnostic or proxy-task accuracy, explanation
quality, evidence consistency, unsupported-claim rate, latency p50 and p95,
failure rate, sample count, and seed. These metrics are not optional because
the paper's central claim concerns reliable decision support rather than text
fluency alone.

\section{Baseline and Ablation Plan}
At least six same-protocol baselines are required before any comparative claim
is allowed: a template-only report, a generic LLM prompt without evidence
grounding, a retrieval-augmented prompt using the same knowledge base, a
structured report without dialogue, a rule-based natural-language explanation,
and a post-hoc feature-importance text summary. A human-written report subset
may be used as an upper-reference comparator if the data protocol permits it.

The ablation suite must remove evidence grounding, retrieval context, dialogue
state tracking, hallucination checking, multi-turn dialogue, and explanation
length controls. The same prompt set, seeds, metrics, and output schema must
be used across all baselines and ablations.

\section{Current Checkpoint Status}
The current code checkpoint passes package-level smoke tests for the local
template path and can run the package-based minimal LLM demo. These checks
prove importability and workflow wiring only. They are not accepted evidence
for manuscript claims because they do not yet emit the required
\texttt{run\_meta.yaml}, \texttt{metrics.json}, latency table, unsupported-claim
table, TOP representative artifacts, or local RTX 4090 provenance.

\begin{table}[!t]
\caption{Current Paper 03 Readiness State}
\label{tab:readiness}
\centering
\begin{tabular}{lll}
\toprule
Gate & Current state & Submission use \\
\midrule
Package import & Passes smoke tests & Wiring only \\
Template demo & Passes local demo & Wiring only \\
IEEE entrypoint & This file & Compile gate \\
Accepted metrics & Missing & Blocker \\
GPU provenance & Missing & Blocker \\
TOP representative & Pending & Blocker \\
SOTA claim & Not allowed & Blocker \\
\bottomrule
\end{tabular}
\end{table}

\section{Limitations and Next Steps}
The main limitations are currently procedural rather than cosmetic. First, the
paper lacks accepted multi-seed evidence packages. Second, hallucination,
retrieval, and latency ablation runners are still missing. Third, TOP recent
work is citation-mapped but not reproduced or represented by accepted local
artifacts. Fourth, current session GPU checks report no visible CUDA devices,
so no accepted local RTX 4090 evidence can be generated from this environment.

The next milestone is to emit accepted evidence packages for the local template
and package-based demo paths, then extend them to same-protocol baselines and
ablations. Only after those artifacts pass the parent artifact and submission
gates should the paper restore quantitative performance, user-study, or SOTA
claims.

\section{Conclusion}
This IEEE Transactions draft establishes a conservative manuscript entrypoint
and evidence contract for the LLM-enhanced explainable fault diagnosis toolkit.
The package is compile-oriented and suitable for readiness review, but it is
not yet a submission-ready paper. Final claims require accepted evidence
packages, same-protocol baselines, ablations, TOP representative artifacts, and
local GPU provenance.

\bibliographystyle{IEEEtran}
\bibliography{references}

\end{document}
